\documentclass[a4paper,12pt]{article}

% ===========================================================================
%  Лабораторная работа 3.4.4 «Петля гистерезиса (статический метод)»
%
%  Сборка:  cd 3.4.4/report && latexmk -pdf main.tex
%  Таблицы и графики читаются из ../data/processed/*.csv, итоговые числа --
%  из ../data/processed/results.tex (макрос \res{...}). Всё это готовит
%  ../scripts/process_data.py. До проведения измерений на месте чисел
%  стоят метки ???.
% ===========================================================================

% ============================================================
% Общий преамбул для отчётов по лабораторным работам (labs_lpr)
% Подключается из report/main.tex командой:
%   % ============================================================
% Общий преамбул для отчётов по лабораторным работам (labs_lpr)
% Подключается из report/main.tex командой:
%   % ============================================================
% Общий преамбул для отчётов по лабораторным работам (labs_lpr)
% Подключается из report/main.tex командой:
%   \input{../../common/latex_template/preamble.tex}
% ============================================================

% Поддержка русского языка и кодировки
\usepackage[T2A]{fontenc}
\usepackage{lmodern}
\usepackage[utf8]{inputenc}
\usepackage[russian]{babel}

% Математические пакеты
\usepackage{amsmath, amssymb}

% Настройка полей
\usepackage{geometry}
\geometry{top=2cm, bottom=2cm, left=2.5cm, right=1.5cm}

% Колонтитулы
\usepackage{fancyhdr}
\usepackage{graphicx}

% Современные таблицы
\usepackage{tabularray}

% Работа с csv файлами
\usepackage{csvsimple-l3}
\usepackage{tikz}

% Графики
\usepackage{pgfplots}
\pgfplotsset{compat=1.18}

% Ссылки
% \usepackage{hyperref} % гипперссылки, нужно оформление вместо красного квадрата
\usepackage[russian]{cleveref} % умные ссылки, нужно переоформление для формул, рисунков и тд (вместо ф-л что-то нужно придумать)

% ============================================================

% Поддержка русского языка и кодировки
\usepackage[T2A]{fontenc}
\usepackage{lmodern}
\usepackage[utf8]{inputenc}
\usepackage[russian]{babel}

% Математические пакеты
\usepackage{amsmath, amssymb}

% Настройка полей
\usepackage{geometry}
\geometry{top=2cm, bottom=2cm, left=2.5cm, right=1.5cm}

% Колонтитулы
\usepackage{fancyhdr}
\usepackage{graphicx}

% Современные таблицы
\usepackage{tabularray}

% Работа с csv файлами
\usepackage{csvsimple-l3}
\usepackage{tikz}

% Графики
\usepackage{pgfplots}
\pgfplotsset{compat=1.18}

% Ссылки
% \usepackage{hyperref} % гипперссылки, нужно оформление вместо красного квадрата
\usepackage[russian]{cleveref} % умные ссылки, нужно переоформление для формул, рисунков и тд (вместо ф-л что-то нужно придумать)

% ============================================================

% Поддержка русского языка и кодировки
\usepackage[T2A]{fontenc}
\usepackage{lmodern}
\usepackage[utf8]{inputenc}
\usepackage[russian]{babel}

% Математические пакеты
\usepackage{amsmath, amssymb}

% Настройка полей
\usepackage{geometry}
\geometry{top=2cm, bottom=2cm, left=2.5cm, right=1.5cm}

% Колонтитулы
\usepackage{fancyhdr}
\usepackage{graphicx}

% Современные таблицы
\usepackage{tabularray}

% Работа с csv файлами
\usepackage{csvsimple-l3}
\usepackage{tikz}

% Графики
\usepackage{pgfplots}
\pgfplotsset{compat=1.18}

% Ссылки
% \usepackage{hyperref} % гипперссылки, нужно оформление вместо красного квадрата
\usepackage[russian]{cleveref} % умные ссылки, нужно переоформление для формул, рисунков и тд (вместо ф-л что-то нужно придумать)


\usepackage{longtable}
\usepackage{caption}
\usepackage{circuitikz}

% Числа на осях: десятичная запятая, без разделителя тысяч
\pgfplotsset{/pgf/number format/.cd, use comma, 1000 sep={}}

% --- Служебные макросы ------------------------------------------------------
\newcommand{\qq}{\textbf{???}}
\newcommand{\fillme}{\textbf{?? НЕОБХОДИМО ЗАПОЛНИТЬ ??}}
\newcommand{\filename}[1]{\texttt{\detokenize{#1}}}
\IfFileExists{../data/processed/results.tex}{% Автоматически создано scripts/process_data.py -- вручную не править.
}{}
\newcommand{\res}[1]{\ifcsname res@#1\endcsname\csname res@#1\endcsname\else\qq\fi}
\newcommand{\cc}{\text{с}}   % индекс «соленоид»

% Колонтитулы (как в остальных работах репозитория)
\pagestyle{fancy}
\fancyhf{}
\fancyhead[L]{%
    \IfFileExists{../../common/latex_template/miptlogo.png}{%
        \includegraphics[height=1cm]{../../common/latex_template/miptlogo.png}%
    }{%
        \textbf{[ЛОГОТИП МФТИ]}%
    }%
}
\fancyhead[R]{\textbf{Емельянов ДА, Байдин ФМ}}
\fancyfoot[C]{\thepage}
\renewcommand{\headrulewidth}{0.4pt}
\setlength{\headheight}{33pt}

\begin{document}

% --- ТИТУЛЬНЫЙ ЛИСТ ---
\begin{titlepage}
\begin{center}
    {\large МОСКОВСКИЙ ФИЗИКО-ТЕХНИЧЕСКИЙ ИНСТИТУТ}
\end{center}
\begin{center}
    {\large Физтех-школа физики и исследований им. Ландау\break ЛФИ}
\end{center}

    \vspace{3cm}

\vspace{0.15cm}
{
\begin{center}
    \huge \bf Лабораторная работа №3.4.4 по курсу Общая физика\\
    \vspace{0.4cm}
    \LARGE Петля гистерезиса (статический метод).
\end{center}
}
\vspace{1.5cm}
\begin{flushright}
Емельянов Дмитрий Алексеевич\\
Байдин Фаддей Михайлович\\
\vspace{0.4cm}
Б02-502
\end{flushright}
\vspace{10cm}
\begin{center}
    \large \qq\ октября 2026, г. Долгопрудный
\end{center}
\end{titlepage}
% ----------------------

\section{Аннотация}
В работе статическим (баллистическим) методом снимаются начальная кривая
намагничивания и предельная петля гистерезиса тороидального образца из
\fillme. Ток в намагничивающей обмотке меняется скачками; изменение
магнитного потока при каждом скачке регистрируется баллистическим
гальванометром, подключённым к измерительной обмотке. Гальванометр
калибруется с помощью длинного пустотелого соленоида.

По петле определяются коэрцитивное поле $H_c$, остаточная индукция $B_r$,
индукция насыщения $B_s$ и намагниченность насыщения $M_s$, по начальной
кривой --- максимальная дифференциальная магнитная проницаемость
$\mu_{\text{диф}}$.

\section{Теоретические сведения}

\subsection{Кривая намагничивания и петля гистерезиса}
Индукция $B$ и напряжённость $H$ в ферромагнетике связаны неоднозначно
(\cref{fig:loop_theory}). При увеличении поля из размагниченного
состояния образец намагничивается по \emph{начальной кривой} $OAC$. Кривая
$OAC$ практически совпадает с $\mu_0 M(H)$, так как $B = \mu_0(H + M)$, а
$\chi \gg 1$. В точке $C$ намагниченность достигает насыщения, и на
участке $CD$
\begin{equation}
    \label{eq:sat}
    B = \mu_0 M_s + \mu_0 H ,
\end{equation}
то есть $CD$ --- прямая с наклоном $\mu_0$, и её экстраполяция к $H = 0$
даёт $\mu_0 M_s$. Индукцию в точке $C$ называют индукцией насыщения $B_s$.

При уменьшении $H$ до нуля система проходит кривую $DCE$; в точке $E$
($H = 0$) сохраняется \emph{остаточная индукция} $B_r$. Чтобы
размагнитить образец (точка $F$), нужно приложить обратное
\emph{коэрцитивное поле} $-H_c$. Замкнутая кривая $CEFC'E'F'C$,
возникающая при циклическом перемагничивании с достижением насыщения, ---
\emph{предельная петля гистерезиса}. Если начать уменьшать поле в
промежуточной точке $A$, получится частная петля $AA'$ меньшей площади.

\begin{figure}[!ht]
    \centering
    \begin{tikzpicture}
        \begin{axis}[
            width=0.62\textwidth, height=0.48\textwidth,
            axis lines=middle, xlabel={$H$}, ylabel={$B$},
            xmin=-4.2, xmax=4.6, ymin=-1.55, ymax=1.65,
            xtick=\empty, ytick=\empty, samples=200, clip=false,
        ]
            \addplot[thick, domain=-3.6:3.6] {tanh((x+0.9)/0.55) + 0.05*x};
            \addplot[thick, domain=-3.6:3.6] {tanh((x-0.9)/0.55) + 0.05*x};
            \addplot[thick, dashed, domain=0:3.6] {tanh(x*x/1.6) + 0.05*x};
            \addplot[densely dotted, domain=-1.3:1.3] {0.86*tanh((x+0.35)/0.45)};
            \addplot[densely dotted, domain=-1.3:1.3] {0.86*tanh((x-0.35)/0.45)};
            \fill (axis cs:0,0.934) circle (1.2pt) (axis cs:-0.9,0) circle (1.2pt)
                  (axis cs:1.3,0.85) circle (1.2pt) (axis cs:-1.3,-0.85) circle (1.2pt)
                  (axis cs:0,-0.934) circle (1.2pt) (axis cs:0.9,0) circle (1.2pt);
            \node[above left]  at (axis cs:3.6,1.18) {$D$};
            \node[above left]  at (axis cs:2.4,1.1) {$C$};
            \node[below right] at (axis cs:-2.4,-1.1) {$C'$};
            \node[below right] at (axis cs:-3.6,-1.18) {$D'$};
            \node[below right] at (axis cs:0,0.95) {$E$};
            \node[above left]  at (axis cs:0,-0.95) {$E'$};
            \node[below left]  at (axis cs:-0.9,0) {$F$};
            \node[above right] at (axis cs:0.9,0) {$F'$};
            \node[right] at (axis cs:1.35,0.78) {$A$};
            \node[left]  at (axis cs:-1.35,-0.78) {$A'$};
            \node[below right] at (axis cs:0,0) {$O$};
            \draw[gray, dashed] (axis cs:0,1.0) -- (axis cs:3.6,1.18);
            \node[left] at (axis cs:0,1.08) {\small $\mu_0 M_s$};
        \end{axis}
    \end{tikzpicture}
    \caption{Петля гистерезиса ферромагнетика (схематично): $OE = B_r$,
             $OF = H_c$}
    \label{fig:loop_theory}
\end{figure}

Наклон кривой намагничивания характеризует дифференциальная магнитная
проницаемость
\begin{equation}
    \label{eq:mu_diff}
    \mu_{\text{диф}} = \frac{1}{\mu_0}\frac{dB}{dH}.
\end{equation}

\subsection{Измерение напряжённости}
На тороидальный сердечник (\cref{fig:toroid}) равномерно намотана
намагничивающая обмотка из $N$ витков. По теореме о циркуляции
напряжённость на средней линии тора
\begin{equation}
    \label{eq:H}
    H = \frac{N}{\pi D}\, I ,
\end{equation}
где $D$ --- средний диаметр тора (поле однородно при $d \ll D$).
Для тора с замкнутым сердечником размагничивающий фактор равен нулю, и
$H$ в образце совпадает с полем $H_0$, создаваемым токами обмотки.

\begin{figure}[!ht]
    \centering
    \begin{tikzpicture}[scale=0.9]
        \draw[thick] (0,0) circle (1.9);
        \draw[thick] (0,0) circle (1.1);
        \foreach \a in {0,15,...,345}
            \draw ({1.0*cos(\a)},{1.0*sin(\a)}) -- ({2.0*cos(\a+4)},{2.0*sin(\a+4)});
        \draw[<->] (-1.5,0) -- (1.5,0) node[midway, above] {$D$};
        \draw[<->] (1.1,-0.35) -- (1.9,-0.35) node[midway, below] {$d$};
        \node at (0,-2.4) {$N$, $N'$};
    \end{tikzpicture}
    \caption{Тороидальный образец с обмотками}
    \label{fig:toroid}
\end{figure}

\subsection{Измерение индукции баллистическим гальванометром}
Поверх намагничивающей обмотки намотана измерительная из $N'$ витков.
При скачке тока $\Delta I$ меняется поле $H$ и поток через измерительную
обмотку, в ней возникает ЭДС
$\mathcal{E} = -SN'\,dB/dt$, $S = \pi d^2/4$. Гальванометр в
баллистическом режиме отклоняется на величину, пропорциональную
протекшему заряду: $\Delta x = \Delta q/b$, где $b$ --- баллистическая
постоянная. Если $R$ --- полное сопротивление измерительной цепи,
\begin{equation}
    \label{eq:dx}
    |\Delta x| = \frac{1}{b}\int I\,dt = \frac{1}{bR}\int\mathcal{E}\,dt
    = \frac{SN'}{bR}\,|\Delta B| .
\end{equation}

Баллистическую постоянную исключают калибровкой: вместо тороида
подключают длинный пустой соленоид ($N_\cc$ витков, длина $l_\cc$,
диаметр $d_\cc$) с короткой измерительной катушкой из $N'_\cc$ витков.
В соленоиде $B_\cc = \mu_0 N_\cc I_\cc / l_\cc$, и при изменении тока на
$\Delta I_\cc$
\begin{equation}
    \label{eq:dx_c}
    |\Delta x_\cc| = \frac{S_\cc N'_\cc}{bR_\cc}|\Delta B_\cc|
    = \frac{\mu_0 S_\cc N'_\cc N_\cc}{bR_\cc l_\cc}|\Delta I_\cc| .
\end{equation}
Деля \cref{eq:dx} на \cref{eq:dx_c}, получаем
\begin{equation}
    \label{eq:dB}
    \Delta B = \mu_0\,\frac{N_\cc N'_\cc}{N'}\,\frac{R}{R_\cc}
               \left(\frac{d_\cc}{d}\right)^2
               \frac{\Delta I_\cc}{l_\cc}\,
               \frac{\Delta x}{\Delta x_\cc}
             \equiv k_B\,\Delta x .
\end{equation}
Если при калибровке уменьшить сопротивление магазина на $R_\cc$
($R'_{\text{М}} = R_{\text{М}} - R_\cc$), сопротивления измерительных цепей совпадают и
$R/R_\cc = 1$.

Индукция в каждой точке петли получается суммированием скачков:
$B_n = B_0 + \sum_{i\le n}\Delta B_i$. Поскольку предельная петля
симметрична относительно начала координат, индукция в начальной точке
каждой ветви равна половине суммы всех скачков по этой ветви:
$B(C) = \frac12\sum_{C\rightarrow C'}|\Delta B_i|$.

\section{Используемое оборудование}
\begin{itemize}
    \item источник питания и генератор тока со ступенчатым тумблером;
    \item тороид с обмотками: $N = \res{N}$, $N' = \res{Nm}$,
          $D = \res{D}$~см, $d = \res{d}$~см; материал: \fillme;
    \item калибровочный соленоид: $N_\cc = \res{Nc}$, $N'_\cc = \res{Ncm}$,
          $l_\cc = \res{lc}$~см, $d_\cc = \res{dc}$~см,
          $R_\cc = \res{Rc}$~Ом;
    \item баллистический гальванометр с осветителем и шкалой,
          $R_0 = \res{Rg}$~Ом;
    \item амперметры $A_1$, $A_2$: \fillme;
    \item магазин сопротивлений $R_{\text{М}}$;
    \item ЛАТР и разделительный трансформатор (для размагничивания);
    \item переключатели П1, П2, ключи К0, К2.
\end{itemize}

\section{Экспериментальная установка и методика измерений}
Схема установки показана на \cref{fig:setup}. Генератор тока, питаемый от
блока питания, позволяет скачками менять ток в намагничивающей обмотке $N$;
ток измеряется амперметром, направление тока меняется переключателем П1.
В цепь измерительной обмотки $N'$ включены гальванометр Г и магазин
сопротивлений $R_{\text{М}}$, которым регулируется чувствительность. Ключ К2
подключает гальванометр только на время измерения, ключ К0 закорачивает
его для быстрой остановки «зайчика», П2 меняет направление тока через
гальванометр. Для калибровки вместо тороида подключается соленоид.
Для размагничивания тороид подключается к ЛАТРу, ток которого плавно
(за $5\div10$~с) уменьшается от тока насыщения до нуля --- образец
проходит 50 петель в секунду всё меньшей площади и приходит в точку
$H = 0$, $B = 0$.

\begin{figure}[!ht]
    \centering
    \begin{circuitikz}[scale=0.95, transform shape]
        \draw (0,0) to[I, l=\small генератор тока] (0,3)
              to[cspst, l=П1] (2,3)
              to[rmeter, t=A] (3.6,3)
              to[short] (4,3)
              to[cute inductor, l_=$N$] (4,0) -- (0,0);
        \draw[very thick] (4.65,0.2) -- (4.65,2.8);
        \draw[very thick] (4.85,0.2) -- (4.85,2.8);
        \node at (4.75,3.35) {\small тороид};
        \draw (5.5,3) to[cute inductor, l=$N'$] (5.5,0);
        \draw (5.5,3) to[R, l=$R_{\text{М}}$] (8,3) to[cspst, l=К2] (9.5,3)
              to[rmeter, t=Г] (9.5,0) -- (5.5,0);
        \draw (11,3) to[cspst, l=К0] (11,0);
        \draw (9.5,3) -- (11,3);
        \draw (9.5,0) -- (11,0);
    \end{circuitikz}
    \caption{Упрощённая схема установки (переключатель П2 и второй
             амперметр не показаны)}
    \label{fig:setup}
\end{figure}

Порядок измерений:
\begin{enumerate}
    \item Подбирается сопротивление $R_{\text{М}}$, при котором наибольший отброс
          (на участках $EF$, $E'F'$) близок к концу шкалы.
    \item \textbf{Предельная петля.} Начиная с максимального тока
          (точка $C$), ток уменьшается щелчками тумблера; для каждого
          положения записывается ток $I$ и отброс $\Delta x$. Вблизи $E$
          фиксируются отбросы и при размыкании, и при замыкании П1 (смена
          направления тока). Далее ток увеличивается до насыщения другого
          знака ($C'$), направление тока через гальванометр меняется П2, и
          петля проходится до точки $C$. Суммы отбросов по двум ветвям
          должны совпадать (расхождение не более 10~\%).
    \item \textbf{Калибровка.} Вместо тороида подключается соленоид,
          $R'_{\text{М}} = R_{\text{М}} - R_\cc$; при размыкании П1 ($\Delta I_\cc = I_{\max}$)
          измеряется отброс $\Delta x_\cc$.
    \item \textbf{Начальная кривая.} Образец размагничивается, затем ток
          скачками увеличивается от нуля до $I_{\max}$ с записью отбросов.
\end{enumerate}
Ток должен меняться строго монотонно; при сбое обход петли начинается
заново от насыщения (для начальной кривой --- после повторного
размагничивания).

\section{Результаты измерений и обработка данных}

\subsection{Параметры установки}
$R_{\text{М}} = \res{RM}$~Ом, $R'_{\text{М}} = \res{RMp}$~Ом, $R_\cc = \res{Rc}$~Ом,
сопротивление гальванометра $R_0 = \res{Rg}$~Ом. Параметры тороида и
соленоида приведены в разделе 3. Коэффициент пересчёта тока в
напряжённость по \cref{eq:H}:
\begin{equation*}
    H = h_I\,I, \qquad h_I = \frac{N}{\pi D} = \res{hI}~\text{м}^{-1}.
\end{equation*}

\subsection{Калибровка гальванометра}
\begin{table}[!ht]
    \centering
    \caption{Калибровка гальванометра на соленоиде}
    \label{tab:calib}
    \csvreader[
        tabularray = {
            colspec = {|c|c|c|},
            hlines, vlines, row{1} = {font=\bfseries},
        },
        table head = {
            $\Delta I_\cc$, мА & $\Delta x_\cc$, дел &
            $\Delta I_\cc/\Delta x_\cc$, мА/дел \\
        },
    ]%
    {../data/processed/calib_processed.csv}%
    {dI_c_mA=\cI, dx_c=\cX, ratio=\cR}%
    {$\cI$ & $\cX$ & $\cR$}
\end{table}

По данным \cref{tab:calib}
$\Delta I_\cc/\Delta x_\cc = (\res{calRatio})$~мА/дел, и по \cref{eq:dB}
цена деления шкалы по индукции
\begin{equation*}
    k_B = (\res{kB})~\text{мТл/дел}, \qquad \varepsilon_{k_B} \approx \res{epsK}~\%.
\end{equation*}
Погрешность $k_B$ складывается из погрешности отсчёта $\Delta x_\cc$,
амперметра и геометрических параметров ($d^2$, $d_\cc^2$, $l_\cc$).

\subsection{Предельная петля гистерезиса}
Результаты измерений и пересчёта приведены в \cref{tab:loop}
(приложение). Сумма отбросов по ветви $C\rightarrow C'$ равна $\res{sumDown}$~дел,
по ветви $C'\rightarrow C$ --- $\res{sumUp}$~дел; расхождение $\res{sumMism}$~\%
(допустимо до 10~\%). Петля и начальная кривая намагничивания построены на
\cref{fig:loop}.

\begin{figure}[!ht]
    \centering
    \begin{tikzpicture}
        \begin{axis}[
            width=0.9\textwidth, height=0.65\textwidth,
            xlabel={$H$, А/м}, ylabel={$B$, Тл},
            grid=both, minor tick num=1,
            max space between ticks=45,
            extra x ticks={0}, extra y ticks={0},
            extra x tick labels={}, extra y tick labels={},
            extra tick style={grid style={black, thin}},
            legend pos=south east, legend cell align=left,
            scaled x ticks=false,
        ]
            \addplot[thick, mark=*, mark size=1.5pt] table[col sep=comma, x=H, y=B]
                {../data/processed/loop_down.csv};
            \addlegendentry{ветвь $C \rightarrow C'$}
            \addplot[thick, mark=square*, mark size=1.5pt] table[col sep=comma, x=H, y=B]
                {../data/processed/loop_up.csv};
            \addlegendentry{ветвь $C' \rightarrow C$}
            \addplot[thick, dashed, mark=triangle*, mark size=1.8pt]
                table[col sep=comma, x=H, y=B]
                {../data/processed/initial_processed.csv};
            \addlegendentry{начальная кривая}
            \ifcsname res@satK\endcsname
                \addplot[thin, domain=0:\csname res@satHmax\endcsname, samples=2]
                    {\csname res@satK\endcsname*x + \csname res@satB\endcsname};
                \addlegendentry{МНК на участке $CD$}
            \fi
        \end{axis}
    \end{tikzpicture}
    \caption{Предельная петля гистерезиса и начальная кривая намагничивания}
    \label{fig:loop}
\end{figure}

Из графика и таблицы:
\begin{itemize}
    \item коэрцитивное поле (среднее по двум ветвям, линейная интерполяция
          между соседними точками)
          $H_c = (\res{Hc})$~А/м;
    \item остаточная индукция (точки с $I = 0$)
          $B_r = (\res{Br})$~Тл;
    \item индукция в точке $C$ при максимальном поле
          $H_{\max} = (\res{Hmax})$~А/м ($I_{\max} = \res{Imax}$~мА):
          $B_s \approx B(C) = (\res{Bmax})$~Тл;
    \item аппроксимация $\res{nSat}$ точек с наибольшим $|H|$ прямой
          \cref{eq:sat} даёт
          $\mu_0 M_s = (\res{muMs})$~Тл, т.е.
          $M_s = (\res{Ms})$~МА/м; наклон участка в единицах $\mu_0$ равен
          $\res{slopeCD}$.
\end{itemize}
\fillme: если наклон участка $CD$ заметно больше $\mu_0$, насыщение при
доступных токах достигнуто не полностью, и $M_s$ --- оценка снизу/сверху.

Погрешность $B_r$ и $B_s$ складывается из погрешности $k_B$ и
накопленной погрешности отсчётов при суммировании скачков
($\Delta B \sim k_B\,\Delta x_{\text{отсч}}\sqrt{n}$); погрешность $H_c$ ---
из погрешности амперметра и шага тока вблизи точки $F$.

\subsection{Начальная кривая намагничивания}
\begin{figure}[!ht]
    \centering
    \begin{tikzpicture}
        \begin{axis}[
            width=0.8\textwidth, height=0.5\textwidth,
            xlabel={$H$, А/м}, ylabel={$\mu_{\text{диф}}$},
            xmin=0, ymin=0, grid=both, minor tick num=1,
            scaled x ticks=false,
        ]
            \addplot[thick, mark=*] table[col sep=comma, x=H, y=mu_diff]
                {../data/processed/mudiff.csv};
        \end{axis}
    \end{tikzpicture}
    \caption{Дифференциальная проницаемость $\mu_{\text{диф}} =
             \Delta B/(\mu_0\Delta H)$ вдоль начальной кривой}
    \label{fig:mudiff}
\end{figure}

Максимальная дифференциальная проницаемость (наибольший наклон начальной
кривой, \cref{fig:mudiff}) достигается при $H \approx \res{HmuMax}$~А/м:
\begin{equation*}
    \mu_{\text{диф}}^{\max} = \res{muMax}.
\end{equation*}
Наибольшее значение $\mu = B/(\mu_0H)$ вдоль начальной кривой
$\mu_{\max} \approx \res{muBHmax}$, начальная проницаемость (по первой
точке) $\mu_{\text{нач}} \approx \res{muInit}$.

\section{Обсуждение результатов}
\begin{itemize}
    \item Для сравнения: техническое (армко) железо --- $H_c \sim
          50\div100$~А/м, $B_s \approx 2{,}1$~Тл
          ($\mu_0 M_s \approx 2{,}15$~Тл), $\mu_{\max} \sim 10^3\div10^4$;
          конструкционная сталь --- $H_c \sim 10^2\div10^3$~А/м.
          \fillme: сравнить полученные $H_c = \res{Hc}$~А/м,
          $B_r = \res{Br}$~Тл, $\mu_0M_s = \res{muMs}$~Тл со справочными
          данными для материала образца.
    \item Совпадение сумм отбросов по ветвям ($\res{sumMism}$~\%)
          показывает, что ни один отброс не пропущен и дрейф нуля
          гальванометра мал.
    \item Систематические погрешности: неоднородность поля по сечению тора
          ($H\propto 1/r$), конечная длина калибровочного соленоида,
          отклонение режима гальванометра от баллистического (при малом
          $R_{\text{М}}$ --- переуспокоение), неполное размагничивание перед снятием
          начальной кривой.
\end{itemize}

\section{Заключение}
Статическим методом исследованы начальная кривая намагничивания и
предельная петля гистерезиса образца из \fillme:
\begin{itemize}
    \item коэрцитивное поле $H_c = (\res{Hc})$~А/м;
    \item остаточная индукция $B_r = (\res{Br})$~Тл;
    \item индукция при максимальном поле $B_s = (\res{Bmax})$~Тл,
          намагниченность насыщения $\mu_0 M_s = (\res{muMs})$~Тл;
    \item максимальная дифференциальная проницаемость
          $\mu_{\text{диф}}^{\max} = \res{muMax}$.
\end{itemize}
\fillme: вывод (магнитомягкий / магнитожёсткий материал, соответствие
справочным данным).

\clearpage
\appendix
\section{Таблица измерений предельной петли}
\begingroup
\small
\csvreader[
    longtable = |c|c|c|c|c|c|,
    table head = {\caption{Измерения предельной петли (ток после щелчка,
            отброс при щелчке, накопленная индукция)}\label{tab:loop}\\
        \hline
        Ветвь & $I$, мА & $H$, А/м & $\Delta x$, дел & $\Delta B$, мТл & $B$, Тл
        \\\hline\endfirsthead
        \hline
        Ветвь & $I$, мА & $H$, А/м & $\Delta x$, дел & $\Delta B$, мТл & $B$, Тл
        \\\hline\endhead
        \hline\endfoot},
    late after line = \\\hline,
]%
{../data/processed/loop_table.csv}%
{branch=\cBr, I_mA=\cI, H=\cH, dx=\cX, dB=\cdB, B=\cB}%
{\cBr & $\cI$ & $\cH$ & $\cX$ & $\cdB$ & $\cB$}
\endgroup

\section{Таблица измерений начальной кривой}
\begingroup
\small
\csvreader[
    longtable = |c|c|c|c|c|,
    table head = {\caption{Измерения начальной кривой намагничивания}
            \label{tab:initial}\\
        \hline
        $I$, мА & $H$, А/м & $\Delta x$, дел & $B$, Тл & $\mu = B/\mu_0H$
        \\\hline\endfirsthead
        \hline
        $I$, мА & $H$, А/м & $\Delta x$, дел & $B$, Тл & $\mu = B/\mu_0H$
        \\\hline\endhead
        \hline\endfoot},
    late after line = \\\hline,
]%
{../data/processed/initial_processed.csv}%
{I_mA=\cI, H=\cH, dx=\cX, B=\cB, mu=\cM}%
{$\cI$ & $\cH$ & $\cX$ & $\cB$ & $\cM$}
\endgroup

\end{document}
